\chapter{Packages, Versions and Migration}
\label{app:packages}

\noindent
What each distribution owns, what version this book was written against, what
moved where in 1.0, and how to migrate a 0.x codebase without guessing.

\chapterstub{
\textbf{Contents.} (A.1) The distribution map: every package used in the book,
what it owns, what depends on it, and its pinned version. (A.2) The version
lines --- which packages move together and which do not, which is the thing that
catches people. (A.3) What moved to \pkg{langchain-classic} and what its
maintenance status means in practice. (A.4) Provider integration packages and the
capability differences that leak through the abstraction (structured output
support, reasoning content, caching, parallel tool calls). (A.5) Checkpointer
backends compared: in-memory, SQLite, Postgres, Redis --- durability, latency,
operational cost. (A.6) The 0.x to 1.x migration, step by step: inventory,
\pkg{langchain-classic} as a bridge, \api{AgentExecutor} to \api{create\_agent},
memory classes to a checkpointer, what LCEL survives. (A.7) The precondition
nobody wants to hear: build the evaluation set before you migrate, or you cannot
tell a regression from a difference.

\textbf{Note.} Versions are macros defined in \code{preamble.tex} and are never
written into a chapter by hand. The CI job in \code{tools/check\_versions.py}
compares them against PyPI on every build.
}
